\documentclass[10pt,a4paper]{amsart}
\usepackage[margin=19mm]{geometry}
\usepackage{amsmath,amssymb,mathtools}
\usepackage[scr=rsfs,cal=euler]{mathalfa}
\usepackage{xfrac}
\usepackage[unicode,hidelinks,bookmarks=false]{hyperref}
\newcommand{\ie}{i.e.\ }
\newcommand{\cf}{cf.\ }
\newcommand{\resp}{respectively\ }
\newcommand{\Subsection}{\subsection}
\providecommand{\nobreakdash}{\nobreak\hbox{-}}
\setlength{\emergencystretch}{2em}
\multlinegap0pt
\usepackage{amssymb}
\newtheorem{thm}{}
\newtheorem{cor}{}
\title{on the discrete convolution of the liouville and m\"obius functions}
\author{Marco Cantarini, Alessandro Gambini, Alessandro Zaccagnini}
\date{Source: arXiv:2603.10241v1 (CC BY 4.0)}
\begin{document}
\maketitle

\noindent\emph{Source and licence.} on the discrete convolution of the liouville and m\"obius functions. Version arXiv:2603.10241v1 (CC BY 4.0), distributed by arXiv under the Creative Commons Attribution 4.0 International licence, http://creativecommons.org/licenses/by/4.0/, as recorded in the arXiv licence metadata for this exact version and shown on the abstract page https://arxiv.org/abs/2603.10241v1. Full licence notice: \texttt{licenses/arxiv-2603.10241-CC-BY-4.0.txt}.\par\smallskip

\noindent\textit{Editorial scope.} Counted unit: the article's thm thm:convol (source0.tex lines 735--751). The article's complete original proof of it is delivered with it (source0.tex lines 753--871, one \texttt{proof} environment of 119 lines). No mathematics is written, repaired or re-derived: the statement and the proof are the article's own lines, in the article's own notation, and no retained line is substituted or re-flowed.  Retained uncounted as the source-local context the proof needs: lines 321--335; lines 681--691.  Nothing was omitted from any retained span, so the package carries no omission record.  No retained line contains a citation, so the wrapper carries no bibliography and no bibliography entry.  No retained line contains a figure environment, an image-inclusion command or drawing code, so no image is delivered and the package declares no asset dependency.  Editorial formatting only, in addition: an AMS \texttt{amsart} preamble that restates the article's own declarations and macros used by the retained lines, namely the environments \texttt{thm}, \texttt{cor} on separate counters, together with the standard packages those lines need.  The retained environments print in this document as Theorem 1, Corollary 1 in order of appearance, whereas the article prints the same items with the numbering of its own full document.  The complete native source of the article as distributed in the arXiv e-print for this exact version is bundled unchanged as \texttt{sources/196123/source0.tex}.
\begin{thm}
\label{thm:double ser}Assume the Riemann hypothesis (RH), let $\rho$
run over the non-trivial zeros of $\zeta(s)$ and let $f(z)$ be a
complex function defined on the set of the
non-trivial zeros and such that
\[
\sum_{0<\gamma<T}\left|\frac{f\left(\rho\right)}{\rho}\right|=o\left(T^{\alpha}\right),\sum_{0<\gamma<T}\left|\frac{f\left(\overline{\rho}\right)}{\overline{\rho}}\right|=o\left(T^{\alpha}\right)
\]
as $T\rightarrow+\infty,$ for every $\alpha>0.$ Then the double
series
\[
\sum_{\rho_{1}}f\left(\rho_{1}\right)\sum_{\rho_{2}}f\left(\rho_{2}\right)\frac{\Gamma\left(\rho_{1}\right)\Gamma\left(\rho_{2}\right)}{\Gamma\left(\rho_{1}+\rho_{2}+1+k\right)}
\]
converges absolutely if $k\in\mathbb{R},$ $k>1/2$.
\end{thm}

\begin{cor}
\label{cor:IMP}Assume RH and SZ. Then for all $T\geq1$ and $x>0$
we have
\[
L(x)=\frac{x^{1/2}}{\zeta\left(\frac{1}{2}\right)}+\sum_{\rho:\left|\gamma\right|<T}\frac{\zeta\left(2\rho\right)x^{\rho}}{\zeta^{\prime}\left(\rho\right)\rho}+R\left(x,T\right)
\]
where
\[
R\left(x,T\right)\ll1+\frac{x\left(\left|\log\left(x\right)\right|+1\right)}{T}+\frac{x-x^{1/4}}{T^{1-\varepsilon}\log\left(x\right)}.
\]
\end{cor}

\begin{thm}
\label{thm:convol}Assume RH and SZ. Then, for any $x>0$ and sufficiently
small $\varepsilon>0$ we have
\[
  C_\lambda(x)
  =
\sum_{n\leq x}S\left(n\right)\left(x-n\right)=\int_{0}^{x}L\left(y\right)L\left(x-y\right)dy
\]
\[
=\frac{x^{2}\pi}{8\zeta\left(\frac{1}{2}\right)^{2}}+\frac{\sqrt{\pi}}{\zeta\left(\frac{1}{2}\right)}\sum_{\rho}\frac{\zeta\left(2\rho\right)\Gamma\left(\rho\right)x^{\rho+3/2}}{\zeta^{\prime}\left(\rho\right)\Gamma\left(\rho+\frac{5}{2}\right)}
\]
\[
+\sum_{\rho_{1}}\frac{\zeta\left(2\rho_{1}\right)}{\zeta^{\prime}\left(\rho_{1}\right)}\sum_{\rho_{2}}\frac{\zeta\left(2\rho_{2}\right)x^{\rho_{1}+\rho_{2}+1}}{\zeta^{\prime}\left(\rho_{2}\right)}\frac{\Gamma\left(\rho_{1}\right)\Gamma\left(\rho_{2}\right)}{\Gamma\left(\rho_{1}+\rho_{2}+2\right)}+O_{\varepsilon}\left(x^{3/2+\varepsilon}+x\right)
\]
where the single series and the double series over the non-trivial
zeros are absolutely convergent.
\end{thm}

\begin{proof}
First of all we note that, using Corollary \ref{cor:IMP}, we have
\[
\int_{0}^{x}L\left(y\right)L\left(x-y\right)dy
\]
\[
=\int_{0}^{x}\left(\frac{y^{1/2}}{\zeta\left(\frac{1}{2}\right)}+\sum_{\rho:\left|\gamma\right|<T}\frac{\zeta\left(2\rho\right)y^{\rho}}{\zeta^{\prime}\left(\rho\right)\rho}\right)\left(\frac{(x-y)^{1/2}}{\zeta\left(\frac{1}{2}\right)}+\sum_{\rho:\left|\gamma\right|<T}\frac{\zeta\left(2\rho\right)(x-y)^{\rho}}{\zeta^{\prime}\left(\rho\right)\rho}\right)dy
\]
\[
+2\int_{0}^{x}L\left(y\right)R\left(x-y,T\right)dy-\int_{0}^{x}R\left(y\right)R\left(x-y\right)dy.
\]
\[
=J_{1}+J_{2}+J_{3}
\]
say. The integration of $J_{1}$ is trivial since all the products
lead, essentially, to Beta functions. So we have
\[
J_{1}=\frac{x^{2}\pi}{8\zeta\left(\frac{1}{2}\right)^{2}}+\frac{\sqrt{\pi}}{\zeta\left(\frac{1}{2}\right)}\sum_{\rho:\left|\gamma\right|<T}\frac{\zeta\left(2\rho\right)\Gamma\left(\rho\right)x^{\rho+3/2}}{\zeta^{\prime}\left(\rho\right)\Gamma\left(\rho+\frac{5}{2}\right)}
\]
\[
+\sum_{\rho_{1}:\left|\gamma_{1}\right|<T}\frac{\zeta\left(2\rho_{1}\right)}{\zeta^{\prime}\left(\rho_{1}\right)}\sum_{\rho_{2}:\left|\gamma_{2}\right|<T}\frac{\zeta\left(2\rho_{2}\right)x^{\rho_{1}+\rho_{2}+1}}{\zeta^{\prime}\left(\rho_{2}\right)}\frac{\Gamma\left(\rho_{1}\right)\Gamma\left(\rho_{2}\right)}{\Gamma\left(\rho_{1}+\rho_{2}+2\right)}.
\]
Now we study $J_{2}$. It is well-known, under RH, that
\[
L\left(x\right)\ll x^{1/2+\varepsilon},\,\forall\varepsilon>0
\]
as $x\rightarrow+\infty$, so we can conclude that for every $\varepsilon>0$
and for every $x>0$ under RH we have the estimation
\[
L\left(x\right)\ll x^{1/2+\varepsilon}+1.
\]

From Corollary \ref{cor:IMP} and from the fact that $$ \frac{x-x^{1/4}}{\log\left(x\right)}\leq x+1$$ for every $x>0,\,x\neq1$, we have
\[
\int_{0}^{x}L\left(y\right)R\left(x-y,T\right)dy
\]
\[
\ll\int_{0}^{x}y^{1/2+\varepsilon}\left[1+\frac{(x-y)\left(\left|\log\left(x-y\right)\right|+1\right)}{T}+\frac{x-y+1}{T^{1-\varepsilon}}\right]dy
\]
\[
+\int_{0}^{x}\left[1+\frac{(x-y)\left(\left|\log\left(x-y\right)\right|+1\right)}{T}+\frac{x-y+1}{T^{1-\varepsilon}}\right]dy
\]
\[
=J_{21}+J_{22}.
\]
\[
\int_{0}^{x}y^{1/2+\varepsilon}dy\ll_{\varepsilon}x^{3/2+\varepsilon},\,\,\frac{1}{T}\int_{0}^{x}y^{1/2+\varepsilon}(x-y)dy=\frac{x^{5/2+\varepsilon}}{T}B\left(\frac{3}{2}+\varepsilon,2\right)
\]
so we have, essentially, just to consider
\[
\frac{1}{T}\int_{0}^{x}y^{1/2+\varepsilon}(x-y)\left|\log(x-y)\right|dy=\frac{1}{T}\int_{0}^{x}\left(x-y\right)^{1/2+\varepsilon}y\left|\log(y)\right|dy
\]
Note that, if $0<v\leq1$, then
\[
v\left|\log\left(v\right)\right|=v\log\left(1/v\right)<1
\]
hence if $0<x\leq1$
\[
\frac{1}{T}\int_{0}^{x}\left(x-y\right)^{1/2+\varepsilon}y\left|\log(y)\right|dy\leq\frac{1}{T}\int_{0}^{x}\left(x-y\right)^{1/2+\varepsilon}dy=\frac{x^{3/2+\varepsilon}}{T}.
\]
If $x>1$, we have
\[
\frac{1}{T}\int_{0}^{1}\left(x-y\right)^{1/2+\varepsilon}y\left|\log(y)\right|dy\ll\frac{x^{1/2+\varepsilon}}{T}
\]
and
\[
\frac{1}{T}\int_{1}^{x}\left(x-y\right)^{1/2+\varepsilon}y\log(y)dy\leq\log\left(x\right)\frac{1}{T}\int_{0}^{x}\left(x-y\right)^{1/2+\varepsilon}ydy
\]
\[
\ll\frac{\log\left(x\right)x^{5/2+\varepsilon}}{T}.
\]
So we can conclude that
\[
J_{21}\ll_{\varepsilon}x^{3/2+\varepsilon}+\frac{\left(\log\left(x\right)x^{5/2+\varepsilon}\right)_{x>1}+x_{0<x\leq1}^{3/2+\varepsilon}}{T}+\frac{x^{3/2+\varepsilon}+x^{5/2+\varepsilon}}{T^{1-\varepsilon}}
\]
and in a similar way
\[
J_{22}\ll x+\frac{x^{2}\log\left(x\right)_{x>1}+1_{0<x\leq1}}{T}+\frac{x+x^{2}}{T^{1-\varepsilon}}.
\]
It remains to consider $J_{3}.$ We have
\[
J_{3}\ll\int_{0}^{x}\left(1+\frac{y\left(\left|\log\left(y\right)\right|+1\right)}{T}+\frac{y+1}{T^{1-\varepsilon}}\right)
\]
\[
\times\left(1+\frac{\left(x-y\right)\left(\left|\log\left(x-y\right)\right|+1\right)}{T}+\frac{x-y+1}{T^{1-\varepsilon}}\right)dy.
\]
If we define
\[
A_{T}\left(y\right)=C_{T}\left(y\right):=1+\frac{1}{T^{1-\varepsilon}}\ll1
\]
\[
B_{T}\left(y\right):=\frac{y\left(\left|\log\left(y\right)\right|+1\right)}{T}+\frac{y}{T^{1-\varepsilon}},D_{T}\left(y\right):=\frac{\left(x-y\right)\left(\left|\log\left(x-y\right)\right|+1\right)}{T}+\frac{x-y}{T^{1-\varepsilon}}
\]
we get
\begin{equation}
J_{3}\ll x+\int_{0}^{x}D_{T}\left(y\right)dy+\int_{0}^{x}B_{T}\left(y\right)dy+\int_{0}^{x}B_{T}\left(y\right)D_{T}\left(y\right)dy\label{eq:eq ABCD}
\end{equation}
and it is very simple to observe that all the integrals in (\ref{eq:eq ABCD})
are convergent for every $x>0$ and they tend to $0$ as $T\rightarrow+\infty$,
hence the main formula follows.

It remains to observe that, actually, the series converges absolutely.
Indeed, for the first sum, by the classical Stirling formula we
have
\[
\sum_{\rho:\gamma>0}\left|\frac{\zeta\left(2\rho\right)\Gamma\left(\rho\right)}{\zeta^{\prime}\left(\rho\right)\Gamma\left(\rho+\frac{5}{2}\right)}\right|\ll\sum_{\rho:\gamma>0}\left|\frac{\zeta\left(2\rho\right)}{\zeta^{\prime}\left(\rho\right)}\right|\frac{1}{\gamma^{5/2}}
\]
and such series is convergent since, by SZ, we have
\[
\sum_{\rho:0<\gamma<T}\left|\frac{\zeta\left(2\rho\right)}{\zeta^{\prime}\left(\rho\right)}\right|\frac{1}{\gamma^{5/2}}\ll\sum_{\rho:0<\gamma<T}\left|\frac{\zeta\left(2\rho\right)}{\rho\zeta^{\prime}\left(\rho\right)}\right|\frac{1}{\gamma^{3/2}}
\]
\[
\ll\frac{\left(\log T\right)^{3/2}\log\left(\log\left(T\right)\right)}{T^{3/2}}+\int_{14}^{T}\left(\left(\log t\right)^{3/2}\log\left(\log\left(t\right)\right)\right)t^{-5/2}dt
\]
and the last integral is convergent if $T\rightarrow+\infty$. 

Clearly, by symmetry, the considered case $\gamma>0$ is enough. For
the double series the convergence follows from Theorem \ref{thm:double ser}.
\end{proof}

\end{document}
